% !TEX root = ../main.tex

\chapter{Related Work}
\label{cha:related}

\sculpt is inserted into a conversation that has been going on for several decades about how to open programming to non-specialist audiences.
That conversation has produced three families of languages worth reviewing separately, since each one partially resolves the problem \sculpt poses.
We name these families \emph{visual and block-based languages}, \emph{tangible languages for computing education}, and \emph{computing as artistic practice} (\emph{creative computing}).
We close the chapter by positioning \sculpt with respect to the three.

\section{Visual and Block-based Languages}
\label{sec:related:visual}

The line of block-based languages begins with \emph{LogoBlocks}~\cite{begel96} in 1996 and consolidates with \emph{Scratch}~\cite{resnick09}, developed at the MIT Media Lab's \emph{Lifelong Kindergarten Group}.
Scratch replaces textual syntax with draggable blocks whose shape defines what fits with what.
This opens an interesting exploration of syntax: in this kind of language, syntax takes on a geometric dimension.
The sets of rules that govern program construction become physical rules.
That idea has been extended by \emph{Blockly}~\cite{fraser15}, a Google library that generates code in target languages (JavaScript, Python, Dart) from block programs, and by \emph{Snap!}~\cite{harvey13}, a reimplementation of Scratch with first-class functions and customisable blocks, aimed at university education.

These proposals drastically reduce syntactic friction for beginners and are now the tool of choice for early computing education.
Their limitations, however, are consistent with the diagnosis from the previous chapter.
The blocks live on a screen, depend on a keyboard (usually) and a mouse, and, by design, do not admit meaningful aesthetic freedom beyond changing background colours.
The interface has been simplified and computation improved, but its material support has not changed.

\section{Tangible Languages for Education}
\label{sec:related:tangible}

A second family moves programming into the physical world.
\emph{AlgoBlock}~\cite{suzuki95} was one of the first systems to propose physical, connectable blocks as a medium for building programs, mainly for collaborative purposes among children.
\emph{Tern}~\cite{horn07} extended the idea with wooden pieces recognised by computer vision, with no need for embedded electronics.
More recently, \emph{KIBO}~\cite{sullivan15} combines wooden blocks with barcodes read by a mobile robot, and \emph{Cubetto}~\cite{cubetto}, commercialised in multiple markets, uses a grid and magnetic tiles to program the movement of a small robot.
\emph{Osmo Coding}~\cite{osmo} follows a hybrid line: the pieces are physical but execution happens on a tablet that observes them through the camera.

These systems amply demonstrate that tangibility reduces cognitive and motivational barriers for student audiences, in line with prior literature on \emph{tangible user interfaces}~\cite{ishii97,marshall07}.
They share, however, three recurring limitations.
First, they are all designed as \emph{educational tools}: their expressive power is intentionally bounded and applied to a specific domain, usually robotics.
This provides a more concrete, manipulative and stimulating programming experience for students, but the resulting systems are still domain-specific languages, or rather, programmable didactic tools.
Second, the pieces are closed.
They come in a shape, colour and size preset by the manufacturer, and the user usually cannot modify them without destroying or altering their functionality.
That freedom, both aesthetic and functional, is absent by design.
The drive to control the user experience is understandable in an educational context, but it limits expressiveness and the range of users who can use these systems.
The freedom to modify pieces matters for users with accessibility needs, who may require pieces of a particular size, material or morphology to manipulate them.
Third, most of these systems come with a price tag that not every educational institution or individual can afford.
In addition to needing specialised hardware (and sometimes software), the cost of the pieces, their maintenance and access can be a barrier to adoption.
This is especially relevant in low-resource contexts, where access to educational tools is crucial to close learning gaps.

\section{Computing as Artistic Practice}
\label{sec:related:art}

A third tradition approaches the problem from a different angle.
It keeps the textual interface of programming but aims it explicitly at artistic production.
\emph{Processing}~\cite{reas07} offers a Java dialect intended for visual artists and designers, and popularised the idea that writing code is a legitimate form of artistic practice.
\emph{Sonic Pi}~\cite{aaron16} builds on the SuperCollider synthesis engine a Ruby superset aimed at musical \emph{live coding}, originally conceived as an educational tool for schools in England.
The practice of \emph{live coding}, both musical and visual, has been systematised as an artistic-computational movement in its own right~\cite{collins03}, with dedicated festivals and communities (\eg \emph{Algorave}).

These initiatives have shown that programming \emph{can} be art and that there is a trained and enthusiastic audience for that vision.
They share, however, a structural limitation with the first family: they do not abandon the typewriter paradigm.
The \emph{live coding} artist still types on a screen in front of an audience (and gets stressed when the code does not compile); the final artistic output is sound, video or image, not the physical body of the program.

\section{Positioning of SCuLPT}
\label{sec:related:positioning}

\sculpt sits at the seldom-travelled intersection of the three preceding families.
It shares with block-based languages the idea that syntax can be a geometric property of the pieces, essential to building programs.
The natural intuitiveness of fitting pieces together is a powerful resource for assimilating complex syntactic concepts and grounds them in a concrete, understandable construction logic: pieces either fit or they do not.
It shares with tangible languages the bet on moving programming off the screen and into three-dimensional space.
Enabling the world to program necessarily implies opening programming to forms of access other than the traditional one.
The aim is to leverage the manipulative and open nature of the pieces to open programming to users with diverse capabilities, resources and accessibility needs.
Finally, it shares with the \emph{creative computing} tradition the conviction that artistic practice is a legitimate vehicle for computation, not a later ornament.
Table~\ref{tab:related} summarises the differences.

\begin{table}[h]
\centering
\small
\caption{Comparison of \sculpt with related families.}
\label{tab:related}
\begin{tabular}{@{}lcccc@{}}
\toprule
                          & Block/visual & Tangible edu. & Creative comp. & \sculpt \\
\midrule
Physical support          & No  & Yes & No  & Yes \\
Turing-complete           & Partial & No & Yes & Yes \\
Aesthetic freedom         & Low & Low & Medium & Very high \\
Target audience           & Students/Novices & Children & Artists & Mixed \\
\bottomrule
\end{tabular}
\end{table}

To the best of our knowledge, \sculpt is the first language to combine \emph{physical tangibility}, \emph{formally proven Turing-completeness} and \emph{radical aesthetic freedom} in a single design.
The radicalness of the aesthetic openness deserves to be underlined: in \sculpt the user can freely modify the shape, size, colour, material and texture of the pieces; the only constrained property is the topology of the connections.
A piece may be a standard PLA cube or an animal figure carved in wood, and the program will work the same as long as the connection points respect the specification.

\endinput
